\documentclass[11pt,a4paper]{article}
\usepackage{amsmath,amssymb,amsthm}
\usepackage{mathtools}
\usepackage{booktabs}
\usepackage{enumitem}
\usepackage{geometry}
\usepackage{xcolor}
\usepackage[pdfencoding=auto,psdextra]{hyperref}

\geometry{margin=2.5cm}

\newcommand{\E}{\mathbb{E}}
\newcommand{\KL}{\mathrm{KL}}
\newcommand{\sg}{\mathrm{sg}}
\newcommand{\xh}{\hat{x}_0}
\newcommand{\real}{\mathrm{real}}
\newcommand{\fake}{\mathrm{fake}}

\newtheorem*{remark*}{Remark}

\title{\textbf{OPD, Approach~B, and TDM: Jacobians for trajectory distribution matching}\\[0.3em]
\large Why swapping MSE for a score difference is not enough, and how truncated
$\partial_\theta x$ recovers reverse-KL along an ODE grid}
\author{Jayce-Ping}
\date{2026-08-02}

\begin{document}
\maketitle

\begin{abstract}
Companion to \texttt{dmd\_opd\_gradient\_relation.tex}. We contrast four objectives used
in XOPD-family distillation: (i)~pointwise OPD field regression; (ii)~endpoint DMD on a
consistency sampler; (iii)~\emph{Approach~B} --- OPD's ODE grid with a score-difference
force on trajectory states; (iv)~Trajectory Distribution Matching (TDM). The critical
invariants for (iii)--(iv) are an independent online fake score, differentiation through
the sampler state ($\partial_\theta x_{t_i}$), and (for TDM) non-overlapping diffusion
intervals. A residual swap that keeps OPD's fixed-$x_t$ Jacobian is \emph{not} reverse KL.
\end{abstract}

\section{Jacobian table}

\begin{center}
\small
\begin{tabular}{@{}llp{5.2cm}@{}}
\toprule
Objective & Differentiates & Match object \\
\midrule
OPD / \texttt{x0\_norm} & $\partial_\theta v_\theta$ at fixed $x_t$ & Teacher field (pointwise) \\
XDMD (endpoint) & $\partial_\theta x0_G$ (one consistency step) & $p_\theta(x_0)$ vs $p_{\real}$ \\
Approach~B (\texttt{xopd\_dm}) & $\partial_\theta x_{t_i}$ (one ODE step) & $p_\theta(x_{t_i})$ vs $p_{\real}$ on OPD grid \\
TDM (\texttt{xtdm}) & $\partial_\theta x_{t_i}$ (one ODE step) & Marginal KL on non-overlapping $[t_i,t_{i+1}]$ \\
\bottomrule
\end{tabular}
\end{center}

\section{Why residual swap $\neq$ KL}

OPD at fixed $x_t$ minimizes $\E\norm{v_\theta(x_t)-v_\psi(x_t)}^2$ (or the $\xh$ form).
Replacing the residual $v_\theta-v_\psi$ by $s_{\fake}-s_{\real}$ while still backpropagating
only through the velocity \emph{output} yields a hybrid with no clean reverse-KL identity:
the score difference is a force in $x$-space, not a field target at a pinned input.

Genuine distribution matching (Diff-Instruct / DMD / TDM) uses
\[
\nabla_\theta\KL(p_\theta\|p_{\real})
\approx \E\big[(s_{\fake}(x)-s_{\real}(x))\,\partial_\theta x\big],
\]
with $s_{\fake}\approx\nabla\log p_\theta$ trained online. Pinning $\fake\leftarrow$teacher
zeros the signal; pinning $\fake\leftarrow$student biases the score and collapses toward
reweighted field matching through $\partial_\theta x$.

\section{Approach~B (ODE grid + score difference)}

Keep the student as a multi-step flow on the OPD inference schedule $\{t_i\}$.
At a randomly chosen index $i^\star$ (broadcast across ranks):
\begin{enumerate}[leftmargin=1.4em]
\item No-grad Euler rollout to $x_{t_{i^\star}}$.
\item One differentiable ODE step producing $x$ (state carrying $\partial_\theta$).
\item Diffuse $x_\tau=(1-\sigma_\tau)x+\sigma_\tau\varepsilon$ with $\tau$ in a band tied to
the selected segment.
\item Evaluate $\xh^{\real},\xh^{\fake}$ under \texttt{no\_grad}; form the DMD2
self-normalized direction
$\mathrm{grad}=(\xh-\xh^{\real})-(\xh-\xh^{\fake})$ (normalized by $\mean|\cdot|$).
\item Stop-grad identity on the \emph{sample}:
$\mathcal{L}=\tfrac12\norm{x-\sg(x-\mathrm{grad})}^2$, so
$\partial\mathcal{L}/\partial x=\mathrm{grad}$.
\end{enumerate}
Memory is $O(1)$ transformer activations for the generator update (truncated BPTT).

\section{TDM (Luo et al., ICCV 2025)}

Paper objective (data-free, deterministic $K$-step generator):
\[
\mathcal{L}(\theta)=\sum_{i=0}^{K-1}\sum_{\tau\in[t_i,t_{i+1}]}\lambda_\tau\,
\KL\big(p_{\theta,\tau|t_i}\,\|\,p_{\phi,\tau}\big).
\]
Intervals $[t_i,t_{i+1}]$ are \emph{non-overlapping} so one fake network can separate
trajectory segments by timestep. Practical surrogate: Pseudo-Huber on
$\norm{x_{t_i}-\sg(\tilde x_{t_i})}$ with
$\tilde x_{t_i}=x_{t_i}+\lambda(s_{\real}-s_{\fake})$.
Engineering (v1): match a \emph{random} segment per generator step (same memory as Approach~B);
full sum over $i$ is optional later.

\section{Shared engineering invariants}

\begin{itemize}[leftmargin=1.4em]
\item Independent online \texttt{fake} LoRA + manual data-parallel optimizer (outside ZeRO).
\item Force acts on sample state $x$, never ``MSE$(v_\theta,v_{\mathrm{teacher}})$ with a score residual''.
\item Autocast weight cache off while swapping teacher weights; DDP bypass for
\texttt{no\_grad} weight swaps when applicable.
\item Collective-safe random segment index (broadcast from rank 0).
\item Eval uses the ODE sampler (multi-step flow), not XDMD's consistency schedule.
\end{itemize}

\section{Relation to capacity}

Under infinite capacity, OPD, Approach~B, and TDM share the teacher field / induced
trajectory as a common fixed point. Under a capacity gap, $L^2$ field projection and
reverse-KL on trajectory marginals diverge: the latter is mode-seeking on the transported
densities, which is why multi-time DM can outperform pure OPD for small students even when
distilling the same deterministic teacher.

\end{document}
